\ejercicio{1,75 puntos}

Considere el siguiente grafo dirigido:

\begin{center}
  \begin{tikzpicture}[-{Latex[length=2.5mm]}, thick, y=0.9cm]
  \node[nodo] (C) at (2, 3.2) {$C$};
  \node[nodo] (A) at (2, 1.2) {$A$};
  \node[nodo] (B) at (0, 0) {$B$};
  \node[nodo] (D) at (4, 0) {$D$};
  \draw (A) -- node[peso] {$2$} (B);
  \draw (A) -- node[peso] {$5$} (C);
  \draw (C) -- node[peso] {$-4$} (B);
  \draw (B) -- node[peso] {$4$} (D);
  \draw (D) -- node[peso] {$1$} (A);
  \draw (D) -- node[peso] {$2$} (C);
\end{tikzpicture}%

\end{center}

\begin{enumerate}[label=(\alph*)]
  \item (0,25 puntos) Ejecute Dijkstra desde $A$. Rompa empates alfabéticamente. Llene una fila por cada extracción, con el vértice extraído y el valor de $d$ justo después de relajar sus arcos.

  \tablatraza{5}{Paso & Extrae & $d[A]$ & $d[B]$ & $d[C]$ & $d[D]$}{1, 2, 3, 4}

  \pagebreak

  \item (0,5 puntos) Ejecute Bellman-Ford desde $A$ relajando los arcos, en cada ronda, en este orden: $B \to D$, $C \to B$, $D \to C$, $D \to A$, $A \to B$, $A \to C$. Llene $d$ al final de cada ronda, incluida la ronda de verificación de ciclos negativos.

  \tablatraza{4}{Ronda & $d[A]$ & $d[B]$ & $d[C]$ & $d[D]$}{1, 2, 3, 4}

  \item (0,25 puntos) Compare los resultados de (a) y (b). ¿Cuál es correcto? ¿En qué paso se equivoca el otro algoritmo y por qué?

  \respuesta[height fill]

  \pagebreak

  \item (0,5 puntos) Ejecute Floyd-Warshall. Las filas son el origen y las columnas el destino. Llene la matriz después de usar cada vértice $k$ como intermedio.

  \begin{center}
    \matrizfw{$D^{(0)}$}{$0$ & $2$ & $5$ & $\infty$}{$\infty$ & $0$ & $\infty$ & $4$}{$\infty$ & $-4$ & $0$ & $\infty$}{$1$ & $\infty$ & $2$ & $0$}
    \hspace{0.5cm}
    \matrizfwvacia{$k = A$}
    \hspace{0.5cm}
    \matrizfwvacia{$k = B$}

    \bigskip
    \matrizfwvacia{$k = C$}
    \hspace{1cm}
    \matrizfwvacia{$k = D$}
  \end{center}

  \item (0,25 puntos) ¿Qué relación hay entre la fila $A$ de la matriz final y el resultado de (b)? Si necesitara las distancias entre todos los pares de vértices de un grafo denso con pesos negativos, ¿usaría Floyd-Warshall o correría Bellman-Ford desde cada vértice? Justifique con la complejidad.

  \respuesta[height fill]
\end{enumerate}
